\section*{Problemas}



\begin{problema}
	Definamos $$p\veebar q \equiv \left( p \vee q \right) \wed \neg\left( p \wed q \right)$$
	Demuestre que 
	\begin{enumerate}
		\item $ \displaystyle
		x \in A\symdif B \iff \left( x \in A \right) \veebar \left( x \in B  \right)
		$ 
		\item $\displaystyle p\veebar q \equiv \left( p \wed \neg q \right) \vee \left( q \wed \neg p \right)$ 
		\item $\displaystyle A \veebar B = \left( A \minus B \right) \cup \left( B \minus A \right)$
	\end{enumerate}
	
\end{problema}

\begin{problema}
	\label{lip:exmp:1.5}
	Supongamos que $\N$ es el conjunto universo. Definamos $A=\set{1,2,3,4},$ $B=\set{3,4,5,6,7},$ $C=\set{2,3,8,9},$ $E=\set{2,4,6,...}.$
	
	Determine:
	\begin{enumerate}
		\item $A \symdif B$ 
		\item $A \symdif C$ 
		\item $B \symdif C$ 
		\item $A \symdif E$
	\end{enumerate}
\end{problema}


\begin{problema}
	\label{lip:exmo:1.6}
	Contruya un sistema de conjuntos fundamentales a partir de tres conjunto $A, B, C.$  
\end{problema}


\begin{problema}
	Encuentre el dual de 
	\[ (\uset \cap A) \cup (B\cap A) = A\]
\end{problema}

\begin{problema}
	Demuestre las siguientes leyes del álgebra de conjuntos:
	\begin{enumerate}
		\item $A \cup A^{C} = \uset$
		\item $A \cap A^{C} = \emptyset$
		\item $(A^{C})^{C} = A$
		\item $A \setminus B = A \cap B^{C}$
		\item $A \symdif B = B \symdif A$ (conmutatividad de la diferencia simétrica)
	\end{enumerate}
\end{problema}

\begin{problema}
	Demuestre las leyes de De Morgan para conjuntos:
	\begin{enumerate}
		\item $(A \cup B)^{C} = A^{C} \cap B^{C}$
		\item $(A \cap B)^{C} = A^{C} \cup B^{C}$
	\end{enumerate}
\end{problema}

\begin{problema}
	Simplifique las siguientes expresiones:
	\begin{enumerate}
		\item $(A \cup B) \cap (A \cup B^{C})$
		\item $(A \cap B) \cup (A \cap B^{C})$
		\item $(A \cup B) \cap A$
		\item $A \cup (A^{C} \cap B)$
	\end{enumerate}
\end{problema}

\begin{problema}
	Sean $A, B, C$ conjuntos. Demuestre que:
	\begin{enumerate}
		\item Si $A \subset B,$ entonces $A \cup C \subset B \cup C$
		\item Si $A \subset B,$ entonces $A \cap C \subset B \cap C$
		\item $A \symdif B = \emptyset \iff A = B$
		\item $A \symdif B = (A \setminus B) \cup (B \setminus A)$
	\end{enumerate}
\end{problema}

\subsection*{Cardinalidad}

\begin{problema}
	En una escuela con 200 estudiantes:
	\begin{itemize}
		\item 120 estudian inglés
		\item 80 estudian francés
		\item 50 estudian alemán
		\item 40 estudian inglés y francés
		\item 30 estudian inglés y alemán
		\item 20 estudian francés y alemán
		\item 10 estudian los tres idiomas
	\end{itemize}
	\begin{enumerate}
		\item ?`Cuántos estudiantes estudian al menos un idioma?
		\item ?`Cuántos no estudian ningún idioma?
		\item ?`Cuántos estudian exactamente un idioma?
	\end{enumerate}
\end{problema}

\begin{problema}
	Sean $A$ y $B$ conjuntos finitos tales que $\card{A} = 20,$ $\card{B} = 15$ y $\card{A \cup B} = 30.$ Calcule:
	\begin{enumerate}
		\item $\card{A \cap B}$
		\item $\card{A \setminus B}$
		\item $\card{B \setminus A}$
		\item $\card{A \symdif B}$
	\end{enumerate}
\end{problema}

\begin{problema}
	Demuestre que si $A$ y $B$ son conjuntos finitos, entonces
	$$
	\card{A \symdif B} = \card{A} + \card{B} - 2\card{A \cap B}.
	$$
\end{problema}

\begin{problema}
	Una encuesta a 150 personas revela que:
	\begin{itemize}
		\item 90 ven televisión
		\item 70 leen libros
		\item 60 escuchan música
		\item 40 ven televisión y leen libros
		\item 35 ven televisión y escuchan música
		\item 30 leen libros y escuchan música
		\item 15 hacen las tres actividades
	\end{itemize}
	\begin{enumerate}
		\item ?`Cuántas personas hacen exactamente dos de estas actividades?
		\item ?`Cuántas hacen solo una?
		\item ?`Cuántas no hacen ninguna?
	\end{enumerate}
\end{problema}